% Generated draft: requires mathematical and source review.
\documentclass[11pt]{article}
\usepackage[margin=25mm]{geometry}
\usepackage{fontspec}
\setmainfont{Latin Modern Roman}
\usepackage{amsmath,amssymb,mathtools,unicode-math}
\setmathfont{Latin Modern Math}
\usepackage{enumitem,needspace,xurl,hyperref}
\setlength{\parindent}{0pt}
\setlength{\parskip}{5pt}
\setlength{\emergencystretch}{3em}
\newcommand{\N}{\mathbb{N}}
\newcommand{\Z}{\mathbb{Z}}
\newcommand{\Q}{\mathbb{Q}}
\newcommand{\R}{\mathbb{R}}
\newcommand{\C}{\mathbb{C}}
\newcommand{\F}{\mathbb{F}}
\newcommand{\A}{\mathbb{A}}
\newcommand{\PP}{\mathbb{P}}
\newcommand{\E}{\mathbb{E}}
\newcommand{\eps}{\varepsilon}
\newcommand{\dd}{\,\mathrm{d}}
\begin{document}
\section*{UnsolvedMath ID: 20002322}
\textbf{Source problem code:} catalogue entry
\subsection{Definitions and mathematical statement}
Use countable ordinary differential fields of characteristic zero in the computable language of rings with one derivation $\delta$. Let $\mathsf{DCF}_0$ be their differentially closed theory. A copy has domain $\omega$ and an atomic diagram in this language. For a fixed countable $K\models\mathsf{DCF}_0$, define its presentation-degree spectrum by
\[
 \operatorname{Spec}_{\mathrm{at}}(K)
   =\{\deg_TD(A):A\cong K,\ \operatorname{dom}(A)=\omega\}.
\]
Computable in a degree $\mathbf d$ means diagram Turing reducible to an oracle of that degree; ``every copy computes $\mathbf d$'' allows the reduction to depend on the copy.

\textbf{Question 20, Harizanov.} Given a Turing degree $\mathbf d$, is there a countable differentially closed field $K$ having a $\mathbf d$-computable copy while every copy computes $\mathbf d$? Equivalently,
\[
 \exists\mathbf e\in\operatorname{Spec}_{\mathrm{at}}(K)\ (\mathbf e\le_T\mathbf d),
 \qquad\forall\mathbf e\in\operatorname{Spec}_{\mathrm{at}}(K)\ (\mathbf d\le_T\mathbf e).
\]
Together these imply $\mathbf d\in\operatorname{Spec}_{\mathrm{at}}(K)$ and that it is the least degree in this spectrum. No uniform functional over all presentations is required.

\textbf{Supplied characteristic-zero answer.} The aid invokes Marker--Miller, Corollary 5.3, in the form
\[
 \forall K\models\mathsf{DCF}_0\text{ countable}\ \forall\mathbf b>\mathbf0\
 \exists\mathbf e\in\operatorname{Spec}_{\mathrm{at}}(K)
                         \quad\mathbf b\not\le_T\mathbf e.
\]
Thus no such spectrum is contained in any nonzero upper cone. The requested existence of \emph{some} $K$ with both clauses is therefore possible exactly for $\mathbf d=\mathbf0$, witnessed by a computable model of $\mathsf{DCF}_0$. This does not assert that every countable differentially closed field has a computable presentation. Positive characteristic or several derivations are outside this supplied result.

\textbf{Supplied countable simultaneous avoidance extension.} The aid further asserts
\[
 \forall K\models\mathsf{DCF}_0\text{ countable}\
 \forall\langle\mathbf b_i:i<\omega\rangle\ (\mathbf b_i>\mathbf0)\
 \exists\mathbf e\in\operatorname{Spec}_{\mathrm{at}}(K)\
 \forall i<\omega\quad\mathbf b_i\not\le_T\mathbf e.
\]
There is one presentation degree avoiding every prescribed nonzero cone, not a separate degree for each cone. This is stronger than the individual avoidance statement; the aid attributes a finite-extension argument to it but gives no full proof or separate reference here. The assertion is existential, with no uniform effective procedure from arbitrary degree representatives.

\textbf{Separate Section 1.12, Andrews.} The appended heading ``Spectrum of totally categorical theories'' defines
\[
 \operatorname{Spec}(T)=\{\mathbf d:\text{some model of }T
                                      \text{ is computable in }\mathbf d\}.
\]
This is an upward oracle-computability spectrum ranging over models of a theory, distinct from exact diagram degrees of copies of one $K$. It is not an assumption of Question 20. The registry cleaned field inserts section-heading text into ``differentially closed''; the original question and supplied interpretation are preserved instead.
\subsection{Short English statement}
Harizanov asks for a differentially closed field with a prescribed least presentation degree. Under the supplied ordinary characteristic-zero result, such a field exists exactly for the computable degree. The aid's stronger extension gives one copy avoiding any countable list of nonzero cones. The adjacent theory-spectrum definition is a different notion.
\subsection{Sources}
\begin{itemize}
\item Local UnsolvedMath registry, permanent ID 20002322, source code catalogue entry; canonical \texttt{aim-logic-notes.json}, notes[97], statement status reconstructed unverified.
\item AIM \emph{Computable stability theory}, Question 20 (Harizanov), followed by Section 1.12 (Andrews): \url{https://aimath.org/pastworkshops/computestabproblems.pdf}. The cleaned field also contains a misplaced Section 1.11/Calvert heading.
\item Marker--Miller, Corollary 5.3, as supplied by the aid; full title, publication details and exact theorem text are absent. The countable simultaneous extension has no separate proving reference in the record.
\item Registry research attempt 1, dated 2026-08-14: full solution, problem status at run solved, rigor ``proved'', novelty confidence low. The explicitly machine-generated, non-peer-reviewed aid requires independent verification before citation.
\end{itemize}
\end{document}
